\section*{Introduction}
\phantomsection\addcontentsline{toc}{section}{Introduction}

TOPORET showed that a population graph and topological features improve DR grading, but
it is not safe enough for screening: its Severe-DR Non-Referral Rate of $3.4\%$ is above
the $2\%$ threshold, because every image is graded, even when the model is unsure. This
chapter presents the main system of the thesis, DOTS (Deep Ordinal Transformer
System)~\cite{belhadj2026dots}, which adds the third axis that TOPORET lacks: ordinal
prediction with guaranteed rank consistency and a calibrated right to abstain.

We first review related work on retinal foundation models, ordinal regression and
abstention in medical AI. We then analyse the residual gaps of TOPORET and state the
hypothesis H3 tested in this chapter. We then present the DOTS architecture, its equations and its training and
inference procedure, followed by the experimental protocol. The results are reported in
four steps: sufficiency (the seven clinical criteria), necessity (can any axis be
removed?), effect sizes and clinical impact, and comparison with the literature. A
discussion of what the results establish, and what they do not, closes the chapter.

\section{Related Work}
\label{sec:ch4-related}

\paragraph{Foundation models for retinal imaging.}
Most DR grading systems are initialised with networks pre-trained on natural images, such
as EfficientNet~\cite{tan2019efficientnet}. RETFound~\cite{zhou2023retfound} changed this
practice: it is a vision transformer pre-trained by masked
autoencoding~\cite{he2022mae} on $1.6$ million unlabelled retinal images, and it improves
several downstream ophthalmic tasks with little labelled data. Such foundation models
address the mismatch between natural-image features and retinal anatomy, but they still
grade each image independently and give no guarantee on the order of the grades.

\paragraph{Ordinal regression for DR grading.}
DR grading is often treated either as plain classification or as regression on the grade
value. Ordinal methods decompose the problem into cumulative binary decisions; among them,
CORAL~\cite{cao2020coral} enforces rank consistency through shared weights and ordered
biases, while CORN~\cite{shi2023corn} models conditional probabilities and relies on the
training loss to keep the grades ordered. To our knowledge, rank consistency has not
previously been guaranteed by construction in a graph-based DR grading system.

\paragraph{Uncertainty and abstention in medical AI.}
Estimating the uncertainty of a deep network, for example with Monte Carlo
dropout~\cite{gal2016dropout} or deep ensembles~\cite{lakshminarayanan2017ensembles}, makes
it possible to refer uncertain cases to an expert. For DR detection, referring the most
uncertain images has been shown to improve the accuracy on the remaining
ones~\cite{leibig2017uncertainty}, and selective classification provides a formal frame for
the trade-off between coverage and error~\cite{geifman2017selective}. Autonomous systems
already cleared by regulators for DR detection~\cite{abramoff2018pivotal} show that
automated screening can reach clinical use, but they do not combine relational context,
vascular topology and abstention. DOTS brings these elements together and evaluates them
against predefined safety criteria.

\section{Motivation and Problem Formulation}
\label{sec:ch4-motivation}

\subsection{Residual Gap of TOPORET}
\label{sec:ch4-mot}

\Cref{fig:dots} gives the complete DOTS architecture that this chapter builds, stage by
stage and axis by axis.

% Figure 4.1 --- complete DOTS architecture
\begin{figure}[!tb]
\centering
\resizebox{\linewidth}{!}{%
\begin{tikzpicture}[
  font=\footnotesize,
  box/.style={draw, rounded corners=2.5pt, align=center, inner sep=2pt,
              minimum height=1.5cm, text width=1.7cm},
  card/.style={draw, rounded corners=2pt, align=center, font=\scriptsize,
               text width=4.3cm, inner sep=4pt},
  arr/.style={-{Latex[length=2mm]}, thick}
]
\def\dx{2.27}
\node[box, fill=black!4, draw=black!55] (in) at (0,0) {\textbf{Input}\\fundus\\image $I_i$};
\node[box, fill=black!6, draw=black!55] (s1) at (\dx,0) {\textbf{S1}\\CLAHE,\\skeleton};
\node[box, fill=axisOf, draw=axisO] (s2) at (2*\dx,0) {\textbf{S2 (O)}\\RETFound\\ViT-L/16};
\node[box, fill=fusionf, draw=fusion] (s4) at (3*\dx,0) {\textbf{S4}\\fusion\\1031-D};
\node[box, fill=axisRf, draw=axisR] (s5) at (4*\dx,0) {\textbf{S5 (R)}\\graph $+$\\GAT, $k{=}15$};
\node[box, fill=axisOf, draw=axisO] (s6) at (5*\dx,0) {\textbf{S6 (O)}\\CORAL $+$\\abstention};
\node[box, fill=okf, draw=okgreen] (out) at (6*\dx,0) {\textbf{Output}\\grade 0--4\\or defer};
\node[box, fill=axisSf, draw=axisS] (s3) at (2.5*\dx,-2.0) {\textbf{S3 (S)}\\$H_0/H_1$\\7-D TDA};
\foreach \a/\b in {in/s1,s1/s2,s2/s4,s4/s5,s5/s6,s6/out} \draw[arr] (\a) -- (\b);
\draw[arr, axisS] (s1.south) |- (s3.west);
\draw[arr, axisS] (s3.east) -| (s4.south);
\node[font=\scriptsize, text=black!65, align=center] at (5.5*\dx,-1.55) {defer if $H_i>0.73$ nats};

\node[card, fill=axisRf, draw=axisR] at (0.9*\dx,-3.75) {\textbf{Axis R closes L1}\\neighbours share signal\\H1: $t_4=3.82$, $p=0.009$};
\node[card, fill=axisSf, draw=axisS] at (3*\dx,-3.75) {\textbf{Axis S closes L2}\\vessel loops $H_0/H_1$\\H2: $F(4,4995)=1308.9$};
\node[card, fill=axisOf, draw=axisO] at (5.1*\dx,-3.75) {\textbf{Axis O closes L3}\\CORAL: $0\%$ rank violations\\Sev-NR $=1.9\%$ $[1.4,\,2.4]$};
\end{tikzpicture}}
\caption[Complete DOTS architecture]{Complete DOTS architecture, read left to right.
Blue: Axis~R (graph, L1); green: Axis~S (topology, L2); purple: Axis~O (ordinal $+$
safety, L3); grey and orange: shared pre-processing and fusion. The skeleton produced by
S1 feeds the topological branch S3, whose 7-D vector is fused with the RETFound
embedding in S4. The three cards name the limitation each axis closes and the result
that confirms it.}
\label{fig:dots}
\end{figure}


TOPORET (\cref{ch:toporet}) confirms H1 and H2: a population graph and topological
features each contribute independently to ordinal grading accuracy. Yet TOPORET is not a
clinically deployable system. At Sev-NR $=3.4\%$, a programme screening $500{,}000$
patients a year would leave about $850$ of its $25{,}000$ severe cases unreferred
(\cref{sec:impact}) -- a rate incompatible with the $2\%$ threshold adopted in this thesis
from national screening practice~\cite{scanlon2017english}.

DOTS (Deep Ordinal Transformer System) is the architectural response to this
residual gap: a system designed, from the outset, around the seven predefined
clinical-oriented criteria presented in \cref{sec:h3def}, rather than around accuracy
maximisation alone. The transition from TOPORET to DOTS is not a single change
but a coordinated modification of three independent subsystems, each addressing
one of three residual gaps identified through systematic error analysis of
TOPORET predictions on the validation set.

\subsection{Three Residual Gaps and Their Solutions}
\label{sec:ch4-gaps}

TOPORET resolves L1 and L2 but leaves three architectural gaps that prevent
clinical deployment. \Cref{fig:gaps} names each gap, the DOTS component
that closes it, and the measured effect. \Cref{sec:gap-a,sec:gap-b,sec:gap-c} take the rows
in order.

% Figure 4.2 --- three residual gaps of TOPORET, closed by DOTS
\begin{figure}[!tb]
\centering
\resizebox{\linewidth}{!}{%
\begin{tikzpicture}[
  font=\footnotesize,
  cell/.style={draw, rounded corners=2.5pt, align=center, inner sep=4pt,
               minimum height=1.55cm, text width=4.0cm},
  arr/.style={-{Latex[length=2mm]}, thick}
]
\def\dx{4.9}\def\dy{2.0}
\node[font=\small\bfseries, text=failred] at (-\dx,0) {Gap left by TOPORET};
\node[font=\small\bfseries, text=navy] at (0,0) {What DOTS adds};
\node[font=\small\bfseries, text=okgreen] at (\dx,0) {Measured effect};

\node[cell, fill=failf, draw=failred] (a1) at (-\dx,-1.2) {\textbf{A -- domain mismatch}\\ImageNet backbone sees\\textures, not the retina};
\node[cell, fill=black!4, draw=black!50] (a2) at (0,-1.2) {RETFound ViT-L/16\\pre-trained on $1.6$M\\retinal images};
\node[cell, fill=okf, draw=okgreen] (a3) at (\dx,-1.2) {\textbf{$+4.1$ pp QWK}\\Gap A closed};

\node[cell, fill=failf, draw=failred] (b1) at (-\dx,-1.2-\dy) {\textbf{B -- ordinal violations}\\CORN: $2.1\%$ of rankings\\are logically impossible};
\node[cell, fill=black!4, draw=black!50] (b2) at (0,-1.2-\dy) {CORAL monotone bias\\(algebraic constraint,\\not a loss penalty)};
\node[cell, fill=okf, draw=okgreen] (b3) at (\dx,-1.2-\dy) {\textbf{$0\%$ violations}\\Gap B closed};

\node[cell, fill=failf, draw=failred] (c1) at (-\dx,-1.2-2*\dy) {\textbf{C -- safety unmet}\\no right to abstain;\\Sev-NR $=3.4\%$};
\node[cell, fill=black!4, draw=black!50] (c2) at (0,-1.2-2*\dy) {dynamic GAT $+$\\entropy abstention\\($\tau_U=0.73$ nats)};
\node[cell, fill=okf, draw=okgreen] (c3) at (\dx,-1.2-2*\dy) {\textbf{Sev-NR $=1.9\%$}\\target met at the point\\estimate (CI up to $2.4\%$)};

\foreach \a/\b in {a1/a2,a2/a3,b1/b2,b2/b3,c1/c2,c2/c3} \draw[arr] (\a) -- (\b);
\end{tikzpicture}}
\caption[The three residual gaps of TOPORET]{The three residual gaps of TOPORET, measured on its ordinal variant under the protocol of this chapter. Read
each row left to right: the problem, the DOTS fix, the measured effect.
\Cref{sec:gap-a,sec:gap-b,sec:gap-c} take the rows in order.}
\label{fig:gaps}
\end{figure}


\subsubsection{Gap A: Domain Mismatch}
\label{sec:gap-a}

EfficientNet-B3, used throughout TOPORET, is pre-trained on ImageNet-1k -- a
corpus of $1.28$ million natural images spanning $1{,}000$ object categories
largely unrelated to medical imaging. The learned filters are optimised for
textures, edges and shapes common in photographs of everyday objects.
Retinal fundus images share almost none of these visual statistics.

RETFound~\cite{zhou2023retfound}, by contrast, is pre-trained by masked
autoencoding (MAE)~\cite{he2022mae} directly on $1.6$ million retinal images. Substituting the
frozen backbone alone, with all other TOPORET components unchanged, yields
QWK $0.880\to 0.921$ (component ablation, \cref{sec:ablation4}), a $+4.1$~pp improvement
attributable entirely to domain-matched pre-training.

\subsubsection{Gap B: Ordinal Rank Violations}
\label{sec:gap-b}

TOPORET uses a CORN ordinal head~\cite{shi2023corn}. At inference time nothing in the
architecture prevents $\hat P(Y\geq 2)>\hat P(Y\geq 1)$, a logically impossible
ordinal violation. Empirically, $2.1\%$ of TOPORET test predictions exhibit
such violations. In a deployed system this can appear as inconsistent
confidence reporting to a clinician, undermining clinical trust and the
IEC~62304 audit trail.

CORAL resolves this through the monotone bias projection of \cref{eq:coral-proj},
an architectural constraint rather than a training-time penalty, yielding
exactly $0\%$ violations by construction (verified on $10{,}000$ test
predictions in \cref{sec:h3a-eval}).

\subsubsection{Gap C: Unmet Safety Threshold}
\label{sec:gap-c}

Even with Gaps~A and~B addressed, a system that must classify every input
will misclassify some genuinely ambiguous cases. The Grade~$2/3$ boundary is
the most clinically ambiguous (inter-rater $\kappa\approx 0.55$--$0.65$).
DOTS addresses Gap~C with two coordinated mechanisms: (i)~a dynamic graph
attention network (\cref{sec:gat}) that refines node embeddings using
per-edge attention recomputed at inference time; and (ii)~an entropy-based
abstention mechanism (\cref{sec:abst}) that defers the
highest-uncertainty cases to human review rather than forcing a potentially
unsafe automated decision.

\subsection{Hypothesis H3: Sufficiency and Necessity}
\label{sec:h3def}

The third hypothesis of the thesis has two parts. (H3a) The three-axis combination satisfies all seven predefined clinical-oriented
criteria simultaneously. (H3b) No proper subset of the three axes does so.
\emph{Pre-specified tests:} (H3a) five-benchmark evaluation against seven predefined
criteria; (H3b) exhaustive enumeration of the $2^3-2=6$ proper non-empty subsets.

\section{Proposed Architecture}
\label{sec:ch4-arch}

\subsection{CORAL Loss with Class-Weighted Ordinal Regression}

The severe class imbalance documented in \cref{app:datasets} (Grade~3 and Grade~4 jointly
constitute under $10\%$ of most benchmarks) motivates class-weighted training.
The DOTS loss is
\begin{equation}
\mathcal{L}_{\mathrm{CORAL}}
= -\,\lambda_{y_i}\sum_{k=1}^{K-1}
\Bigl[
 y_i^{(k)}\log\sigma(g_i^{(k)})
+(1-y_i^{(k)})\log\bigl(1-\sigma(g_i^{(k)})\bigr)
\Bigr],
\label{eq:coral}
\end{equation}
with $g_i^{(k)}=w^\top z'_i+b_k$ and the monotone bias projection
$b_k\leftarrow\min(b_k,b_{k-1}-\epsilon)$ (\cref{eq:coral-proj}). The loss of
each training image is multiplied by the weight $\lambda_{y_i}$ of its true grade; the
class weights $(\lambda_0,\ldots,\lambda_4)=(1.0,2.5,3.0,8.0,6.0)$ were selected by
grid search on the validation folds, minimising Sev-NR while keeping the best macro-F1,
with no class weight exceeding $10\times$ the baseline. \Cref{tab:weights} compares the selected weights with uniform and
inverse-frequency weighting: the capped scheme gives the best macro-F1 and is the only
one that brings Sev-NR below $2\%$.

\begin{table}[H]
\centering
\caption[Class-weight sensitivity]{Class-weight sensitivity: validation macro-F1 and Sev-NR for three
weighting schemes. The selected weights balance class-wise sensitivity against
overall calibration.}
\label{tab:weights}
\small
\begin{tabular}{llcc}
\toprule
Scheme & Weights $(\lambda_0,\ldots,\lambda_4)$ & Macro-F1 & Sev-NR \\
\midrule
Uniform & $(1,1,1,1,1)$ & $0.812$ & $4.8\%$ \\
Inverse-frequency & $(1.0,9.4,3.2,28.1,16.8)$ & $0.834$ & $2.6\%$ \\
Selected (capped) & $(1.0,2.5,3.0,8.0,6.0)$ & $0.847$ & $1.9\%$ \\
\bottomrule
\end{tabular}
\tabnote{Row~3 is the $5\times 5$ grid-search
winner.}
\end{table}

\subsection{Dynamic Graph Attention}
\label{sec:gat}

TOPORET uses a static population graph. DOTS replaces this with a dynamic graph
attention layer that recomputes per-edge attention at each inference call:
\begin{equation}
z'_i=\mathrm{GAT}(z_i,G)=\sum_{j\in\mathcal N(i)}\alpha_{ij}Wz_j,
\label{eq:gat}
\end{equation}
\begin{equation}
\alpha_{ij}
=\frac{\exp\bigl(\mathrm{LeakyReLU}(a^\top[Wz_i\|Wz_j])\bigr)}
{\sum_{k\in\mathcal N(i)}\exp\bigl(\mathrm{LeakyReLU}(a^\top[Wz_i\|Wz_k])\bigr)}.
\label{eq:att}
\end{equation}
Unlike the static TOPORET graph, $\alpha_{ij}$ is patient-specific: two query
patients sharing neighbour $j$ may attend to $j$ with different weights.
``Neighbour $j$ contributed $\alpha_{ij}=0.34$ to this prediction'' is therefore
a well-defined inference-time quantity (EU AI Act Article~13).

\subsection{Entropy-Based Safety Abstention}
\label{sec:abst}

All parameters affecting the final test-set decision rule were selected on
training and validation folds only. The held-out test set was accessed exactly
once, for the final evaluation. Patient $i$ is deferred when
\begin{equation}
H_i=-\sum_c p_{ic}\log(p_{ic}+\epsilon)> \tau_U=0.73~\mathrm{nats}.
\label{eq:ent}
\end{equation}
The threshold is the smallest $\tau_U$ on the validation sweep
$\tau_U\in[0.3,1.2]$ that satisfies Sev-NR $<2\%$ subject to a coverage floor
of $90\%$. \Cref{tab:tau} shows the sweep: lower thresholds defer too many images
(coverage below $90\%$), higher ones let too many severe cases through, and
$\tau_U=0.73$ is the only tested value that satisfies both constraints.

\begin{table}[H]
\centering
\caption[Threshold sweep for entropy abstention]{Threshold sweep for entropy abstention. $\tau_U=0.73$ is the minimal
threshold satisfying both the Sev-NR and coverage constraints.}
\label{tab:tau}
\small
\begin{tabular}{cccc}
\toprule
$\tau_U$ (nats) & Coverage & Sev-NR & Constraint? \\
\midrule
$0.40$ & $78.1\%$ & $1.1\%$ & Fails coverage floor \\
$0.55$ & $86.4\%$ & $1.5\%$ & Fails coverage floor \\
$0.73$ & $93.8\%$ & $1.9\%$ & Satisfies both \\
$0.90$ & $97.2\%$ & $2.6\%$ & Fails Sev-NR \\
$1.10$ & $99.5\%$ & $3.9\%$ & Fails Sev-NR \\
\bottomrule
\end{tabular}
\tabnote{Coverage rises, and safety falls,
monotonically with $\tau_U$.}
\end{table}

\Cref{fig:tradeoff} plots the same values as a trade-off curve.

% Coverage--safety trade-off of the abstention threshold (data of Table tab:tau)
\begin{figure}[H]
\centering
\begin{tikzpicture}
\begin{axis}[width=0.78\linewidth, height=6.2cm,
  xlabel={Coverage (\% of images graded automatically)}, ylabel={Sev-NR (\%)},
  xmin=75, xmax=101, ymin=0.5, ymax=4.3,
  tick label style={font=\footnotesize}, label style={font=\small},
  grid=major, grid style={black!10}]
  \fill[okf, opacity=0.8] (axis cs:90,0.5) rectangle (axis cs:101,2.0);
  \node[font=\scriptsize, text=okgreen, anchor=south east] at (axis cs:100.8,0.55) {both criteria met};
  \draw[failred, densely dashed, thick] (axis cs:75,2.0) -- (axis cs:101,2.0)
     node[pos=0.03, anchor=south west, font=\scriptsize] {Sev-NR $=2\%$};
  \draw[axisR, densely dashed, thick] (axis cs:90,0.5) -- (axis cs:90,4.3)
     node[pos=0.97, anchor=north east, font=\scriptsize, align=right] {coverage\\$=90\%$};
  \addplot[very thick, axisO, mark=*, mark options={fill=white}]
    coordinates {(78.1,1.1) (86.4,1.5) (93.8,1.9) (97.2,2.6) (99.5,3.9)};
  \node[font=\scriptsize, anchor=south east] at (axis cs:78.1,1.15) {$\tau_U{=}0.40$};
  \node[font=\scriptsize, anchor=south east] at (axis cs:86.4,1.55) {$0.55$};
  \node[font=\scriptsize\bfseries, anchor=north west] at (axis cs:93.8,1.85) {$0.73$ (selected)};
  \node[font=\scriptsize, anchor=north west] at (axis cs:97.2,2.55) {$0.90$};
  \node[font=\scriptsize, anchor=east] at (axis cs:99.2,3.9) {$1.10$};
\end{axis}
\end{tikzpicture}
\caption[Coverage--safety trade-off of the abstention threshold]{Coverage--safety
trade-off of the abstention threshold $\tau_U$ (values of \cref{tab:tau}). Raising the
threshold grades more images automatically but lets more severe cases through; the green
region contains the settings that satisfy both the coverage floor ($90\%$) and the safety
threshold ($2\%$). Only $\tau_U=0.73$ lies inside it.}
\label{fig:tradeoff}
\end{figure}


\subsection{Design Rationale}
\label{sec:ch4-rationale}

\paragraph{A frozen foundation model.}
RETFound is used as a frozen feature extractor rather than fine-tuned end to end. Fine-tuning
a ViT-Large model on each fold would multiply the computational cost and increase the risk
of overfitting the smaller datasets; keeping it frozen also makes the S2 stage a fixed,
versioned component, which simplifies verification (\cref{sec:iec62304}). The gain of Gap~A
is therefore obtained without any retinal-specific training of the backbone within this
thesis.

\paragraph{CORAL rather than CORN.}
TOPORET used CORN, which models conditional probabilities and usually trains well, but
whose rank consistency depends on the learned weights. In a safety-critical setting, an
inconsistent ranking -- for example a higher probability of ``at least Grade~4'' than of
``at least Grade~3'' -- is not only an error but an output that cannot be explained to a
clinician. CORAL removes this failure mode by construction, at the cost of a shared weight
vector for all thresholds; the ablation shows that this constraint does not reduce accuracy
(\cref{sec:ablation4}).

\paragraph{Dynamic rather than static attention.}
In TOPORET, the weight of an edge is fixed once the graph is built. In DOTS, the GAT layer
recomputes the importance of each neighbour for each query patient (\cref{eq:att}). This
allows the model to rely on different neighbours for different patients, and it produces,
for every prediction, a ranked list of the neighbours that influenced it, which serves
directly as an explanation.

\paragraph{Entropy as the abstention criterion.}
Several uncertainty estimates could decide when to abstain: Monte Carlo
dropout~\cite{gal2016dropout} or deep ensembles~\cite{lakshminarayanan2017ensembles}
capture epistemic uncertainty better, but require several forward passes per image, which
multiplies the latency. The predictive entropy of the CORAL output is available in a single
pass, is well defined for ordinal outputs, and, because DOTS is well calibrated
(ECE $=0.040$), tracks the actual probability of error closely enough to meet the safety
criterion. Combining it with ensemble-based uncertainty is a natural extension.

\subsection{Training and Inference Procedure}

\Cref{alg:dots} gives the complete training and inference procedure. Training follows the
TOPORET pipeline with the three changes above; at inference time, the only additional step
is the entropy test that decides whether the grade is returned or the case is deferred. The
first part of the algorithm (lines~1--6) is run once per cross-validation fold: it computes
the fused features of every training image, builds the population graph, trains the
attention layer and the ordinal head, and selects the abstention threshold on the
validation data. The second part (lines~7--14) is run for every new patient: it retrieves the
neighbours of the new image in the existing graph, refines its representation with the
attention weights, computes the five grade probabilities and their entropy, and either
returns a grade or defers the case to a human grader. The two parts share the same
feature extractors, so that a new patient is processed exactly like the training images.

\begin{algorithm}[H]
\caption{DOTS: training and safe inference with abstention.}
\label{alg:dots}
\DontPrintSemicolon
\KwIn{training set $\{(I_j,y_j)\}$; new image $I_i$; class weights $\lambda_0,\ldots,\lambda_4$; threshold $\tau_U=0.73$}
\KwOut{grade $\hat y_i$, or \texttt{DEFER} with the explanation of \cref{sec:audit}}
\BlankLine
\tcp{Training}
\ForEach{training image $I_j$}{
  $z_j \leftarrow [\,\mathrm{RETFound}(\mathrm{CLAHE}(I_j)) \,\Vert\, x_j^{\mathrm{TDA}}\,]$\tcp*{stages S1--S4}
}
Build the dual-similarity $k$-NN graph $G$ (\cref{eq:dualsim}, $k=15$)\tcp*{stage S5}
Train the dynamic GAT layer and the CORAL head with the weighted loss of \cref{eq:coral}\;
After each update, enforce $b_k \leftarrow \min(b_k,b_{k-1}-\epsilon)$\tcp*{\cref{eq:coral-proj}}
Select $\tau_U$ on the validation folds (\cref{tab:tau})\;
\BlankLine
\tcp{Inference for a new patient}
Compute $z_i$; retrieve its $k$ neighbours in $G$\;
$z'_i \leftarrow \mathrm{GAT}(z_i,\mathcal{N}(i))$ with attention weights $\alpha_{ij}$\tcp*{\cref{eq:gat,eq:att}}
$p_{i,c} \leftarrow$ grade probabilities from the CORAL head\;
$H_i \leftarrow -\sum_c p_{i,c}\log p_{i,c}$\tcp*{\cref{eq:ent}}
\eIf{$H_i > \tau_U$}{\Return{\texttt{DEFER} to a human grader}}{\Return{$\hat y_i \leftarrow \arg\max_c p_{i,c}$}}
\end{algorithm}

Training uses the settings of \cref{app:hyper}: the Adam optimiser with a learning rate of
$10^{-3}$, mini-batches of $32$ target
nodes, $100$ epochs and five folds stratified by grade, with the
same random seed as TOPORET. The RETFound and topological features are computed once and
cached, so that each epoch only updates the graph layer and the CORAL head; this keeps the
cost of each training run moderate (\cref{tab:compute}). At inference time, the cost
of grading a new patient does not depend on the size of the cohort: the features of the new
image are computed once, its neighbours are retrieved with an approximate nearest-neighbour
query, and the GAT layer aggregates a fixed number of sampled neighbours. The complete
procedure takes $31.4$~ms per image (\cref{tab:latency}), dominated by the topological
features, and the abstention test itself adds a negligible cost, since the entropy is
computed from the five grade probabilities that the model outputs anyway. The same code path is used for evaluation and for deployment, so the measured performance corresponds to the behaviour of the system that would be installed.

\subsection{Regulatory-Oriented Design}
\label{sec:regdesign}

The architectural choices above also address key regulatory obligations \emph{by
construction}, not by retrofit: the typed stages support IEC~62304 unit verification, the
per-stage typed exceptions support ISO~14971 risk management, the GAT weights
$\alpha_{ij}$ provide the per-prediction transparency of EU AI Act Article~13, and entropy
abstention implements the human oversight of Article~14. \Cref{ch:reg} develops this
mapping in full (\cref{tab:regmap}).

The same design choices also simplify the management of changes after deployment. Because
the RETFound backbone and the topological extractor are frozen and versioned, and because
every stage has a typed contract, a modification can be confined to one stage and verified
by that stage's unit tests, without re-validating the whole pipeline. DOTS is therefore
intended to be deployed as a \emph{locked} algorithm: the weights are fixed at release, and
any retraining produces a new version that goes through the same verification and
evaluation protocol before it replaces the previous one. This is the model expected by
current medical-device regulations for software whose behaviour must remain predictable
once in clinical use.

\section{Experimental Protocol}
\label{sec:ch4-protocol}

\subsection{Datasets, Baselines and Parameters}

DOTS is evaluated on the five public benchmarks of \cref{app:datasets}: Kaggle EyePACS,
APTOS~2019, IDRiD, Messidor-2 and DDR, acquired in four countries with different cameras.
The held-out test set of the H3a evaluation contains $10{,}000$ images from APTOS~2019 and
Kaggle EyePACS; the other benchmarks are used to measure performance across acquisition
conditions.

Four reference systems are compared with DOTS: EfficientNet-B3~\cite{tan2019efficientnet},
the most common convolutional baseline; RETFound~\cite{zhou2023retfound}, the strongest
retinal foundation model; TOPORET (\cref{ch:toporet}), which resolves L1 and L2 only; and
the six reduced systems of the necessity analysis (\cref{sec:h3b}), which remove one or two
of the three axes. All systems share the same data splits, pre-processing and number of
training epochs.

The graph parameters are inherited from TOPORET ($\alpha=0.65$, $k=15$, $L=2$ layers,
$S=25$ sampled neighbours, hidden dimension $256$). The two parameters specific to DOTS are
the class weights of the CORAL loss, $\lambda=(1.0,2.5,3.0,8.0,6.0)$ for Grades~0--4
(\cref{tab:weights}), and the abstention threshold $\tau_U=0.73$~nats (\cref{tab:tau}).
Both were selected on the validation folds only. The complete list of hyperparameters is
given in \cref{app:hyper}.

\subsection{Evaluation Protocol and Contamination Controls}
\label{sec:protocol}

The strength of the DOTS results invites scrutiny of the evaluation protocol.
\Cref{tab:protocol} lists each potential source of test-set contamination and the control
used to block it.

\begin{table}[H]
\centering
\caption[Evaluation-protocol controls]{Evaluation-protocol controls. Each potential confound is listed with
the control used to block it.}
\label{tab:protocol}
\footnotesize
\begin{tabular}{P{4.4cm}P{8.4cm}}
\toprule
Potential confound & Control implemented \\
\midrule
Graph construction leakage & Population graph built from training-fold embeddings only; test nodes are inductive. \\
Threshold $\tau_U$ optimisation & Grid search on the validation fold only; test set accessed once. \\
Class weights $\lambda_0,\ldots,\lambda_4$ selection & Validation-fold Sev-NR minimisation; frozen before test. \\
$\alpha=0.65$ (dual metric) & 5-fold CV on training data only. \\
$k=15$ (neighbours) & Validation-fold QWK maximisation. \\
Multi-seed stability & Means over 3 seeds; variance within $\pm 0.003$ QWK. \\
Benchmark selection & Five datasets selected a priori; none excluded after seeing results. \\
\bottomrule
\end{tabular}
\end{table}

Graph-based models need particular care here. Because each prediction uses the
neighbours of a patient, a test image that was present in the graph during training, or a
threshold tuned on the test set, would leak information into the evaluation and inflate
every metric. The controls below were designed so that the test images influence neither
the graph used for training nor any selected parameter; each row of the table names one
such risk and the rule that blocks it.

All hyperparameters affecting the final decision rule were selected on training and
validation data, and the five test sets were accessed exactly once, after all design
decisions were finalised. The seven evaluation criteria and their thresholds were fixed at
the same time, before any test result was seen: QWK $>0.90$, AUC $>0.97$, ECE $<0.05$,
Sev-NR $<2\%$, PDR-NR $<1.5\%$, coverage $>90\%$ and $0\%$ rank violations. Each criterion
is reported with a bootstrap 95\% confidence interval, but the pre-specified decision rule
uses the point estimate (\cref{app:stats}).

\section{Results}
\label{sec:ch4-results}

\subsection{Sufficiency: The Seven Criteria (H3a)}
\label{sec:h3a-eval}

DOTS is evaluated against seven pre-specified criteria on the combined
APTOS $+$ Kaggle EyePACS test set ($N=10{,}000$); \cref{tab:h3a} lists each criterion,
its threshold and the result. H3a is satisfied at the
point-estimate level, but statistical compliance with the predefined safety
threshold is not established (95\% CI upper bound $2.4\%$ exceeds $2\%$).
The CORAL violation rate is identically zero across $10{,}000$ test
predictions, confirming the algebraic guarantee of \cref{eq:coral-proj}.

\begin{table}[H]
\centering
\caption[H3a: seven pre-specified clinical-oriented criteria]{H3a: seven pre-specified clinical-oriented criteria
($N=10{,}000$ held-out predictions). Pass $=$ the point estimate meets the
threshold.}
\label{tab:h3a}
\small
\begin{tabular}{clclc}
\toprule
\# & Criterion & Threshold & DOTS [95\% CI] & Status \\
\midrule
1 & Quadratic Weighted Kappa & $>0.90$ & $0.951$ $[0.941,\,0.957]$ & Pass \\
2 & AUC (macro one-vs-rest) & $>0.97$ & $0.988$ $[0.982,\,0.994]$ & Pass \\
3 & Expected Calibration Error & $<0.05$ & $0.040$ $[0.034,\,0.046]$ & Pass \\
4 & Severe-DR Non-Referral Rate & $<2.0\%$ & $1.9\%$ $[1.4\%,\,2.4\%]$ & Pass \\
5 & PDR Non-Referral Rate & $<1.5\%$ & $0.8\%$ $[0.4\%,\,1.2\%]$ & Pass \\
6 & Coverage (after abstention) & $>90\%$ & $93.8\%$ $[91.2\%,\,96.4\%]$ & Pass \\
7 & CORAL rank violations & $=0\%$ & $0.0\%$ $[0.0\%,\,0.0\%]$ & Pass \\
\bottomrule
\end{tabular}
\tabnote{The Sev-NR 95\% CI upper bound is $2.4\%$,
so statistical compliance with the $2\%$ line is not established.}
\end{table}

The seven criteria are not arbitrary: QWK and AUC are the standard ordinal and
discriminative metrics of the DR literature; ECE operationalises calibration
for triage; Sev-NR and PDR-NR are taken from published national screening
guidelines; coverage stops abstention from trivially buying safety; rank
violations are the architectural guarantee of CORAL.

\subsection{Cross-Benchmark Performance}
\label{sec:bench}

\Cref{tab:bench} reports the QWK of DOTS and of three reference systems on the five
public benchmarks~\cite{kaggle2015eyepacs,aptos2019,porwal2018idrid,
decenciere2014messidor,li2019ddr}. DOTS has the highest QWK on four of them; on IDRiD,
the smallest dataset ($N=516$), RETFound remains higher ($0.901$ against $0.878$), a gap
discussed among the limitations in the General Conclusion.

\begin{table}[H]
\centering
\caption[Quadratic weighted kappa on each of the five public benchmarks]{Quadratic weighted kappa on each of the five public benchmarks.}
\label{tab:bench}
\small
\begin{tabular}{lccccc}
\toprule
System & EyePACS & APTOS & IDRiD & Messidor-2 & DDR \\
\midrule
EfficientNet-B3 & $0.840$ & $0.851$ & $0.832$ & $0.847$ & $0.838$ \\
RETFound & $0.920$ & $0.928$ & $0.901$ & $0.915$ & $0.912$ \\
TOPORET & $0.880$ & $0.887$ & $0.878$ & $0.900$ & $0.871$ \\
DOTS & $\mathbf{0.951}$ & $\mathbf{0.953}$ & $0.878$ & $\mathbf{0.948}$ & $\mathbf{0.951}$ \\
\bottomrule
\end{tabular}
\end{table}

\Cref{fig:bench} shows the same results graphically. The ranking of the systems is the same on all benchmarks except IDRiD, which suggests that the gains of DOTS do not depend on a particular camera or population.

% QWK per benchmark (data of Table tab:bench)
\begin{figure}[H]
\centering
\begin{tikzpicture}
\begin{axis}[ybar, width=\linewidth, height=6.2cm, bar width=7pt,
  ymin=0.80, ymax=0.97, ylabel={QWK},
  symbolic x coords={EyePACS,APTOS,IDRiD,Messidor-2,DDR}, xtick=data,
  enlarge x limits=0.14,
  legend image code/.code={\draw[#1] (0cm,-0.1cm) rectangle (0.3cm,0.12cm);},
  legend style={at={(0.5,1.02)}, anchor=south, legend columns=4, font=\footnotesize, draw=none},
  tick label style={font=\footnotesize}, label style={font=\small},
  ymajorgrids, grid style={black!10},
  yticklabel style={/pgf/number format/fixed,/pgf/number format/precision=2}]
\addplot[fill=black!25, draw=black!60] coordinates {(EyePACS,0.840) (APTOS,0.851) (IDRiD,0.832) (Messidor-2,0.847) (DDR,0.838)};
\addplot[fill=fusionf, draw=fusion] coordinates {(EyePACS,0.920) (APTOS,0.928) (IDRiD,0.901) (Messidor-2,0.915) (DDR,0.912)};
\addplot[fill=axisRf, draw=axisR] coordinates {(EyePACS,0.880) (APTOS,0.887) (IDRiD,0.878) (Messidor-2,0.900) (DDR,0.871)};
\addplot[fill=axisOf, draw=axisO] coordinates {(EyePACS,0.951) (APTOS,0.953) (IDRiD,0.878) (Messidor-2,0.948) (DDR,0.951)};
\legend{EfficientNet-B3, RETFound, TOPORET (ordinal variant), DOTS}
\draw[axisR, densely dashed] (axis cs:EyePACS,0.90) -- (axis cs:DDR,0.90);
\end{axis}
\end{tikzpicture}
\caption[QWK of the four systems on the five public benchmarks]{QWK of the four systems on
the five public benchmarks (values of \cref{tab:bench}); the dashed line is the
$0.90$ criterion. DOTS is the best system on four benchmarks; on IDRiD, the smallest
dataset, RETFound remains ahead.}
\label{fig:bench}
\end{figure}


\subsection{Component Ablation Study}
\label{sec:ablation4}

To isolate the contribution of each architectural change between TOPORET and
DOTS, a sequential ablation introduces Gaps~A, B and~C one at a time (\cref{tab:ablate}). Each
addition improves both QWK and Sev-NR, and no single addition crosses the $2\%$ line on
its own: the threshold is reached only when the last component is added.

\begin{table}[H]
\centering
\caption[Sequential ablation from TOPORET to DOTS]{Sequential ablation from TOPORET to DOTS on Kaggle EyePACS
($N=88{,}702$, 5-fold mean).}
\label{tab:ablate}
\small
\begin{tabular}{lccc}
\toprule
Configuration & QWK & Sev-NR & Violations \\
\midrule
TOPORET (baseline) & $0.880$ & $3.4\%$ & $2.1\%$ \\
$+$ RETFound backbone (Gap A) & $0.921$ & $3.1\%$ & $1.9\%$ \\
$+$ CORAL ordinal head (Gap B) & $0.934$ & $2.6\%$ & $0.0\%$ \\
$+$ Dynamic GAT (partial Gap C) & $0.943$ & $2.2\%$ & $0.0\%$ \\
$+$ Entropy abstention (full Gap C) $=$ DOTS & $0.951$ & $1.9\%$ & $0.0\%$ \\
\bottomrule
\end{tabular}
\tabnote{Each row adds one Gap resolution. No single addition crosses the
$2\%$ safety line on its own.}
\end{table}

\subsection{Failure-Mode Analysis}
\label{sec:fail}

\subsubsection{Confusion Matrix}

\Cref{tab:confusion} gives the DOTS confusion matrix on the $9{,}380$ test images that were
not deferred (coverage $93.8\%$ of the $10{,}000$ held-out predictions).

\begin{table}[!htb]
\centering
\caption[DOTS confusion matrix on accepted predictions]{DOTS confusion matrix on accepted predictions ($N=9{,}380$). Rows: true grade;
columns: predicted grade; each row sums to $100\%$.}
\label{tab:confusion}
\footnotesize
\begin{tabular}{@{}lccccc@{}}
\toprule
True $\backslash$ Predicted & Grade 0 & Grade 1 & Grade 2 & Grade 3 & Grade 4 \\
\midrule
Grade 0 & \textbf{97.3\%} & 2.5\% & 0.2\% & 0.0\% & 0.0\% \\
Grade 1 & 4.1\% & \textbf{84.2\%} & 11.4\% & 0.3\% & 0.0\% \\
Grade 2 & 0.3\% & 8.7\% & \textbf{82.9\%} & 7.9\% & 0.2\% \\
Grade 3 & 0.0\% & 0.2\% & 11.3\% & \textbf{78.4\%} & 10.1\% \\
Grade 4 & 0.0\% & 0.0\% & 0.6\% & 8.2\% & \textbf{91.2\%} \\
\bottomrule
\end{tabular}
\end{table}

The dominant
error is adjacent-grade confusion (Grade~$k\leftrightarrow k\pm1$), consistent with the
known difficulty of the ICDR scale at grade boundaries, where inter-rater agreement among
ophthalmologists is only moderate ($\kappa\approx0.55$--$0.65$ at the Grade~2/3
boundary). There are no clinically dangerous multi-grade jumps: no Grade~0 image is
classified as Grade~3 or~4, and no Grade~4 image as Grade~0 or~1. This is the empirical
counterpart of the rank-consistency guarantee of CORAL.

\subsubsection{Abstention by True Grade}

\Cref{tab:abstention} decomposes the $6.2\%$ abstention rate by true grade. Deferral
concentrates on genuinely ambiguous cases: $73\%$ of abstained images lie at a grade
boundary, where expert agreement is lowest. Grade~0 has the lowest rate overall ($2.1\%$), and among the diseased grades Grade~4 has
the lowest ($3.1\%$), as expected for a grade with no upper boundary and visually
distinctive neovascularisation.
The high rate at Grade~1 ($14.2\%$) reflects the clinical reality that the Mild/Moderate
boundary is among the most contested in practice.

\begin{table}[H]
\centering
\caption[Abstention rate by true grade]{Abstention rate by true grade. Deferral concentrates on ambiguous grade
boundaries.}
\label{tab:abstention}
\small
\begin{tabular}{@{}lrcl@{}}
\toprule
True grade & $N$ (test) & Abstention rate & Share of deferred cases at a boundary \\
\midrule
Grade 0 & 5{,}234 & 2.1\% & 62\% (G0/G1) \\
Grade 1 & 1{,}021 & 14.2\% & 81\% (G1/G2) \\
Grade 2 & 2{,}489 & 11.6\% & 74\% (G2/G3) \\
Grade 3 & 612 & 9.4\% & 88\% (G3/G4) \\
Grade 4 & 644 & 3.1\% & 41\% (no upper boundary) \\
\midrule
Overall & 10{,}000 & 6.2\% & 73\% \\
\bottomrule
\end{tabular}
\tabnote{A deferred case is counted ``at a boundary'' when its two highest predicted grade
probabilities are adjacent grades and sum to more than $0.85$.}
\end{table}

\subsubsection{Per-Grade Calibration}

\Cref{tab:ece} decomposes the overall ECE of $0.040$ by true grade. Calibration is worst
at Grade~3 and Grade~1, which lie on the contested Grade~2/3 and Mild/Moderate boundaries,
and best at Grade~0 (the majority class) and Grade~4 (visually distinctive). The grades
with high abstention rates are broadly those with high calibration error: entropy
abstention and ECE are two views of the same phenomenon, boundary ambiguity.

\begin{table}[H]
\centering
\caption[Expected Calibration Error by true grade]{Expected Calibration Error and mean confidence of correct predictions, by true
grade.}
\label{tab:ece}
\small
\begin{tabular}{@{}lcc@{}}
\toprule
Grade & Grade-conditional ECE & Mean confidence (correct predictions) \\
\midrule
Grade 0 & 0.021 & 0.96 \\
Grade 1 & 0.058 & 0.81 \\
Grade 2 & 0.052 & 0.84 \\
Grade 3 & 0.061 & 0.79 \\
Grade 4 & 0.034 & 0.91 \\
\midrule
Overall (micro-averaged) & \textbf{0.040} & -- \\
\bottomrule
\end{tabular}
\end{table}

\section{Necessity Analysis (H3b)}
\label{sec:h3b}

DOTS combines three ingredients, summarised in \cref{tab:axes5}: the relational
graph~(R), the vascular topology~(S), and the ordinal-safety axis~(O: RETFound, CORAL and
the right to abstain).

\begin{table}[H]
\centering
\caption[The three axes of DOTS tested by H3b]{The three axes of DOTS tested by H3b.}
\label{tab:axes5}
\small
\begin{tabular}{@{}c l l@{}}
\toprule
Axis & Plain name & What it does \\
\midrule
R & Relational graph & similar patients share information (dual-similarity GraphSAGE) \\
S & Vascular topology & $H_0/H_1$ persistence of the vessel skeleton (7-D vector) \\
O & Ordinal $+$ safety & RETFound, CORAL ranking and entropy abstention \\
\bottomrule
\end{tabular}
\end{table}

Hypothesis H3b asks one concrete question:

\begin{quote}
\emph{If any of these ingredients is removed, can the remaining system still keep Sev-NR
below $2\%$ under the same test protocol?}
\end{quote}

There are six ways to drop something: the three single axes $\{R\}$, $\{S\}$, $\{O\}$ and
the three pairs $\{R,S\}$, $\{R,O\}$, $\{S,O\}$. The full triple $\{R,S,O\}$ is DOTS
itself. All six reduced systems were trained with every other setting held as in DOTS
and scored on the same seven criteria as \cref{tab:h3a}.
\Cref{fig:h3b} summarises the outcome: none of the six reduced systems meets the safety
rule. The figure also shows that the reduced systems fail by very different margins, which the next subsections examine one by one.

% Figure 5.1 — centred H3b diagram
\begin{figure}[H]
\centering
\begin{tikzpicture}[
  font=\small,
  fail/.style={draw=failred, fill=failf, rounded corners=3pt, align=center,
               inner sep=5pt, minimum width=2.7cm, minimum height=1.45cm},
  pass/.style={draw=okgreen, fill=okf, rounded corners=3pt, align=center,
               inner sep=6pt, minimum width=4.4cm, minimum height=2.7cm}
]

% singles
\node[fail] (R)  at (-4.3, 2.35) {\textbf{\{R\} only}\\Sev-NR $=4.1\%$\\fails the $2\%$ rule};
\node[fail] (S)  at (-4.3, 0.45) {\textbf{\{S\} only}\\Sev-NR $=4.0\%$\\fails the $2\%$ rule};
\node[fail] (O)  at (-4.3,-1.45) {\textbf{\{O\} only}\\Sev-NR $=3.6\%$\\fails the $2\%$ rule};

% pairs
\node[fail] (RS) at ( 4.3, 2.35) {\textbf{\{R,S\}}\\Sev-NR $=3.4\%$\\fails the $2\%$ rule};
\node[fail] (RO) at ( 4.3, 0.45) {\textbf{\{R,O\}}\\Sev-NR $=2.2\%$ *\\fails the $2\%$ rule};
\node[fail] (SO) at ( 4.3,-1.45) {\textbf{\{S,O\}}\\Sev-NR $=2.4\%$\\fails the $2\%$ rule};

% DOTS
\node[pass] (D) at (0,0.45) {\textbf{DOTS \{R, S, O\}}\\[2pt]Sev-NR $=\mathbf{1.9\%}$\\95\% CI $[1.4\%,\ 2.4\%]$\\[2pt]{\scriptsize point estimate passes}\\{\scriptsize CI upper bound still $>2\%$}};

\foreach \x in {R,S,O}
  \draw[-{Latex}, failred, thick] (\x.east) -- (D.west);
\foreach \x in {RS,RO,SO}
  \draw[-{Latex}, failred, thick] (\x.west) -- (D.east);

\node[font=\scriptsize\color{failred}] at (0,-2.55) {predefined safety line: Sev-NR $=2.0\%$};
\draw[failred, densely dashed, thick] (-6.1,-2.78) -- (6.1,-2.78);
\node[font=\scriptsize\color{black!60}, text width=12.2cm, align=center] at (0,-3.25)
  {* Closest failure: \{R,O\} at $2.2\%$. Bootstrap 95\% CI on (DOTS $-$ \{R,O\}) is $[-0.54,-0.08]$ pp and excludes $0$.};
\end{tikzpicture}
\caption[H3b: all six proper subsets fail the 2\% rule]{H3b in one picture: within the investigated architectural family, the six ways of dropping an axis all fail Sev-NR $<2\%$ (red boxes); only the full three-axis system (green) reaches $1.9\%$, at the point estimate only. Each arrow compares one reduced system with DOTS.}
\label{fig:h3b}
\end{figure}


\subsection{Implementation of Each Reduced System}

\Cref{tab:subimpl} specifies what was trained for each reduced system. Every component
that is not listed is removed; everything else (data, folds, optimiser, seeds and number
of epochs) is identical to DOTS.

\begin{table}[H]
\centering
\caption[Implementation of each reduced system]{Implementation of each reduced system. Only the listed components are kept.}
\label{tab:subimpl}
\small
\begin{tabular}{@{}cl P{10.2cm}@{}}
\toprule
\# & System & Implementation \\
\midrule
1 & $\{R\}$ & EfficientNet-B3 $+$ visual GraphSAGE. No TDA, no CORAL, no abstention. \\
2 & $\{S\}$ & EfficientNet-B3 $+$ TDA concatenated. No graph, no CORAL, no abstention. \\
3 & $\{O\}$ & EfficientNet-B3 $+$ CORAL. No graph, no TDA, no abstention. \\
4 & $\{R,S\}$ & TOPORET: dual-similarity graph $+$ TDA $+$ CORN (not CORAL). No abstention. \\
5 & $\{R,O\}$ & RETFound $+$ visual GraphSAGE $+$ CORAL $+$ abstention. No TDA. \\
6 & $\{S,O\}$ & RETFound $+$ TDA $+$ CORAL $+$ abstention. No graph. \\
7 & $\{R,S,O\}$ & DOTS (all three axes). \\
\bottomrule
\end{tabular}
\tabnote{Systems 3, 5 and 6 contain the ordinal axis. System 3 keeps only its CORAL head
on the ImageNet backbone; systems 5 and 6 keep the complete axis.}
\end{table}

\subsection{Results of the Exhaustive Enumeration}

\Cref{tab:h3b} reports the Sev-NR of each reduced system on the $10{,}000$-image test
set used for H3a.

\begin{table}[H]
\centering
\caption[H3b: Sev-NR of every reduced system versus DOTS]{Sev-NR of every reduced system versus DOTS. Fail: Sev-NR $\geq 2\%$.}
\label{tab:h3b}
\small
\begin{tabular}{@{}cllc@{}}
\toprule
\# & System & Sev-NR & Verdict \\
\midrule
1 & $\{R\}$ & 4.1\% & Fail \\
2 & $\{S\}$ & 4.0\% & Fail \\
3 & $\{O\}$ & 3.6\% & Fail \\
4 & $\{R,S\}$ (TOPORET) & 3.4\% & Fail \\
5 & $\{R,O\}$ & 2.2\% & Fail (closest) \\
6 & $\{S,O\}$ & 2.4\% & Fail \\
\midrule
7 & $\{R,S,O\}$ (DOTS) & \textbf{1.9\%} & Pass (point estimate) \\
\bottomrule
\end{tabular}
\end{table}

\paragraph{H3b confirmed.}
All six proper subsets fail Sev-NR $<2\%$ (\cref{tab:h3b}). The closest failure is
$\{R,O\}$ at $2.2\%$: graph and ordinal safety, but no topology. It lies $0.3$~pp above
the line, and the bootstrap 95\% interval of the difference with DOTS is
$[-0.54,\,-0.08]$~pp, which excludes zero. Removing topology alone is therefore enough to
miss the safety rule, even when the other two axes are at their strongest.

\subsection{Full Seven-Criterion Breakdown by Subset}

Sev-NR is the binding constraint for every failing subset, but the six other criteria
were also evaluated, to show which criteria each subset fails and by how much
(\cref{tab:sevencrit}).

\begin{table}[H]
\centering
\caption[Full seven-criterion results for the six proper subsets and DOTS]{Full seven-criterion results for the six proper subsets and DOTS
(\cmark{} pass, \xmark{} fail).}
\label{tab:sevencrit}
\footnotesize
\setlength{\tabcolsep}{2.4pt}
\begin{tabular}{@{}lccccccc@{}}
\toprule
Subset & \makecell{QWK\\$>0.90$} & \makecell{AUC\\$>0.97$} & \makecell{ECE\\$<0.05$} & \makecell{Sev-NR\\$<2\%$} & \makecell{PDR-NR\\$<1.5\%$} & \makecell{Cov.\\$>90\%$} & \makecell{Viol.\\$=0\%$} \\
\midrule
$\{R\}$ & 0.852 \xmark & 0.961 \xmark & 0.071 \xmark & 4.1\% \xmark & 2.1\% \xmark & 100\% \cmark & 3.8\% \xmark \\
$\{S\}$ & 0.849 \xmark & 0.958 \xmark & 0.068 \xmark & 4.0\% \xmark & 2.0\% \xmark & 100\% \cmark & 4.2\% \xmark \\
$\{O\}$ & 0.864 \xmark & 0.967 \xmark & 0.052 \xmark & 3.6\% \xmark & 1.7\% \xmark & 100\% \cmark & 0.0\% \cmark \\
$\{R,S\}$ & 0.880 \xmark & 0.972 \cmark & 0.058 \xmark & 3.4\% \xmark & 1.6\% \xmark & 100\% \cmark & 2.1\% \xmark \\
$\{R,O\}$ & 0.921 \cmark & 0.979 \cmark & 0.046 \cmark & 2.2\% \xmark & 1.1\% \cmark & 95.1\% \cmark & 0.0\% \cmark \\
$\{S,O\}$ & 0.918 \cmark & 0.977 \cmark & 0.048 \cmark & 2.4\% \xmark & 1.2\% \cmark & 94.6\% \cmark & 0.0\% \cmark \\
\midrule
DOTS & 0.951 \cmark & 0.988 \cmark & 0.040 \cmark & 1.9\% \cmark & 0.8\% \cmark & 93.8\% \cmark & 0.0\% \cmark \\
\bottomrule
\end{tabular}
\end{table}

The breakdown reveals a structural pattern. Subsets without the ordinal axis ($\{R\}$,
$\{S\}$, $\{R,S\}$) fail on calibration and rank consistency as well as on Sev-NR. The
two subsets that contain the complete ordinal-safety axis but lack one of the other two
($\{R,O\}$, $\{S,O\}$) satisfy six of the seven criteria and fail only on Sev-NR. This is
consistent with the role of each axis: the ordinal-safety axis targets calibration and
rank consistency, while the relational and structural axes improve discrimination at the
Grade~2/3 boundary, the main driver of residual Sev-NR.

\paragraph{Sensitivity analysis of the borderline case $\{R,O\}$.}
The most critical comparison is between $\{R,O\}$ (Sev-NR $=2.2\%$) and DOTS
($1.9\%$), a difference of only $0.3$~pp. Four analyses address whether this margin
supports the necessity of the structural axis.
\begin{enumerate}
\item \textbf{Bootstrap CI:} $\Delta(\mathrm{DOTS}-\{R,O\})=-0.31$~pp, 95\% CI
$[-0.54,\,-0.08]$~pp. The interval excludes zero.
\item \textbf{Threshold sensitivity:} if the threshold were relaxed to $2.5\%$ (a
plausible alternative for some programmes), $\{R,O\}$ would pass. The necessity claim is
therefore contingent on the $2\%$ threshold, which is a policy choice, not a
mathematical law.
\item \textbf{Largest benchmark:} on Kaggle EyePACS alone ($N=88{,}702$), $\{R,O\}$
reaches Sev-NR $=2.18\%$, again above the threshold.
\item \textbf{Scope:} H3b is established constructively within the evaluated
architectural family and specification. It does not claim that no conceivable two-axis
system could satisfy the criterion (\cref{sec:scope}).
\end{enumerate}

\subsection{Scope of the Necessity Claim}
\label{sec:scope}

H3b is an exhaustive check of six engineered systems, not a statistical test over a
sample of architectures. It is constructive and exhaustive within the specification used:
the same five benchmarks, the same seven criteria, the same architecture family (CNN or
ViT backbone $+$ optional graph $+$ optional TDA $+$ optional CORAL $+$ optional
abstention). It does not prove that no conceivable architecture lacking one of the axes
could meet the criteria; that would require a proof over the space of all architectures,
which is neither attempted nor claimed. The result reads: among the best-tuned
implementations of each proper subset achievable in this family, none satisfies all
seven criteria, while the full combination does. This is the standard scope of a
constructive necessity argument in applied machine learning, strengthened here by
exhaustive rather than partial enumeration.

\section{Effect Sizes and Clinical Impact}
\label{sec:ch4-impact}

\subsection{Effect Sizes}
\label{sec:effects}

\Cref{tab:effects} lists the effect size of each primary comparison. All exceed the
conventional ``large'' thresholds (Cohen's $d>0.8$; $\eta^2>0.14$~\cite{cohen1988}).
The comparison with EfficientNet-B3 reaches $d=8.2$, which Cohen's framework would call
extreme. Such a value between two trained systems reflects the cumulative effect of five
architectural changes (\cref{sec:progress}) on retrospective data; it is a property of
this specific comparison, not a general claim about achievable effect sizes, and it
should not be extrapolated to prospective settings before the Phase~2 reader study. Effect sizes are nevertheless useful for planning that study, since they indicate the size of the differences that it will need to detect. A smaller prospective sample would suffice for the largest effects, but the Sev-NR criterion, which concerns rare severe cases, will drive the required sample size.

\begin{table}[H]
\centering
\caption[Effect sizes for the primary comparisons]{Effect sizes for the primary comparisons.}
\label{tab:effects}
\small
\begin{tabular}{@{}llcl@{}}
\toprule
Comparison & Metric & Effect size & Interpretation \\
\midrule
DOTS vs.\ EffNet-B3 & QWK $+11.1$~pp ($0.840\to0.951$) & $d=8.2$ & Extreme \\
DOTS vs.\ RETFound & QWK $+3.1$~pp ($0.920\to0.951$) & $d=1.4$ & Very large \\
Population graph (H1) & QWK $+4.0$~pp & $d_z=1.71$ & Very large \\
Vascular topology (H2) & ANOVA & $\eta^2=0.51$ & Large \\
Necessity (H3b) & Constructive enumeration & $6/6$ subsets fail & -- \\
\bottomrule
\end{tabular}
\end{table}

\subsection{Clinical Impact Projection}
\label{sec:impact}

A projection makes the practical stakes concrete. Sev-NR is the proportion of severe
cases (Grades~3--4) that are \emph{not} referred, so the number of missed cases is
\begin{equation}
\text{missed severe cases per year} \;=\; N\times\pi_{\mathrm{sev}}\times\text{Sev-NR},
\label{eq:impact}
\end{equation}
where $N$ is the number of patients screened per year and $\pi_{\mathrm{sev}}$ the
proportion of them with severe disease. For illustration, \cref{tab:impact} applies \cref{eq:impact} to a
national programme with $N=500{,}000$ and an assumed $\pi_{\mathrm{sev}}=5\%$
($25{,}000$ severe cases), including proliferative disease in $1\%$ of the screened
population ($5{,}000$ cases). These prevalences are assumptions, not measurements; the
relative reduction does not depend on them.

\begin{table}[H]
\centering
\caption[Clinical impact projection for 500,000 screened patients per year]{Clinical impact projection for a programme screening $500{,}000$ patients per
year (assumed prevalence: $5\%$ severe, including $1\%$ proliferative).}
\label{tab:impact}
\small
\begin{tabular}{@{}lccc@{}}
\toprule
Quantity & EfficientNet-B3 & DOTS & Difference \\
\midrule
Sev-NR & 4.3\% & 1.9\% & $-2.4$~pp \\
PDR-NR & 2.0\% & 0.8\% & $-1.2$~pp \\
Missed severe cases (Grades 3--4) per year & 1{,}075 & 475 & $-600$ ($-56\%$) \\
\quad of which proliferative (Grade 4) & 100 & 40 & $-60$ ($-60\%$) \\
Images deferred to a human grader per year & 0 & 31{,}000 & -- \\
\bottomrule
\end{tabular}
\end{table}

Under these assumptions, moving from EfficientNet-B3 to DOTS avoids about $600$ missed
sight-threatening cases per year, a relative reduction of $56\%$ that holds for any
prevalence. The cost is that $6.2\%$ of images (about $31{,}000$ per year) are deferred
to human graders; these are, by construction, the most ambiguous cases. The projection
counts only directly prevented missed referrals, not the downstream value of earlier
treatment.

\paragraph{A caveat on causal interpretation.}
The projection assumes that the retrospective Sev-NR improvement transfers to prospective
deployment at the same magnitude. This is exactly what Phases~2--3 of the translation
roadmap (\cref{tab:roadmap}) are designed to test; it must not be treated as an
established clinical outcome before that validation.

\subsection{Synthesis: Sufficiency, Necessity and Magnitude}
\label{sec:synthesis}

This chapter completes a three-part evidential structure that began in \cref{ch:dots}.
H3a (\cref{sec:h3a-eval}) established sufficiency: the full architecture clears every
threshold at the point estimate. H3b (\cref{sec:h3b}) established necessity within the
evaluated family: no proper subset does. The effect-size and clinical-impact analyses add
a third dimension that sufficiency and necessity alone do not capture: magnitude. A
necessity result could in principle be technically true yet practically marginal, the
best proper subset missing the threshold by a negligible amount. This is not the case
here: the best proper subset, $\{R,O\}$, misses by $0.3$~pp in absolute terms, a relative
Sev-NR increase of about $16\%$ over DOTS, or about $75$ additional missed severe cases
per year in the projection of \cref{tab:impact}.

\section{Comparison with the Literature}
\label{sec:ch4-lit}

\subsection{Comparison with Concurrent Work}
\label{sec:concurrent}

Several 2024--2026 systems pursue related but distinct strategies.
Ensembles improve raw accuracy ($+1$--$2$~pp QWK) but do not address L1, L2 or
L3. Self-supervised retinal pre-training, used alone, addresses only Gap~A.
Post-hoc calibration improves ECE without the algebraic rank-consistency of
CORAL. DOTS's claim is not that any single component is novel, but that this
specific combination, applied to this seven-criterion specification, is both
architecturally coherent and empirically necessary (\cref{ch:dots}, H3b).

\subsection{Progressive Contribution Analysis}
\label{sec:progress}

\Cref{fig:progress} traces the full progression from the 2019
baseline to DOTS, showing the cumulative impact of each contribution on accuracy and
safety. The improvement is monotone in both QWK and Sev-NR: no component degrades the
others. The safety threshold is crossed only at the final step (entropy abstention), and
only when all earlier components are in place. This is the experimental counterpart of
the constructive necessity established by H3b.

The two metrics are not driven by the same components. The largest gain in accuracy comes
from the retinal foundation model ($+0.041$ QWK), whereas the largest single reduction of
the non-referral rate comes from the dual-similarity graph with topological features
($-0.7$~pp, from $4.1\%$ to $3.4\%$), followed by the CORAL head ($-0.5$~pp). Abstention adds
the smallest QWK gain of all steps ($+0.008$) but is the step that crosses the $2\%$
threshold. A design driven by accuracy alone would therefore have stopped after the
foundation model and missed the safety criterion; a design driven by safety needs all
three axes.

% Figure 5.2 --- DOTS ablation on Messidor-2: QWK (bars) and non-referral rates (lines)
\begin{figure}[!htbp]
\centering
\pgfplotsset{
  progaxis/.style={
    width=0.9\textwidth, height=6.4cm,
    xmin=0.4, xmax=5.6, xtick={1,...,5},
    xticklabels={{EfficientNet-B3\\baseline},{CNN $+$\\GraphSAGE},{CNN $+$\\CORAL only},{CNN $+$\\TTA only},{DOTS\\complete}},
    xticklabel style={align=center, font=\scriptsize},
    tick label style={font=\footnotesize}, label style={font=\small},
    axis x line*=bottom,
  }}
\begin{tikzpicture}
\begin{axis}[progaxis,
    ymin=0.86, ymax=0.96, ylabel={QWK (bars)}, axis y line*=left,
    ytick={0.86,0.88,0.90,0.92,0.94},
    yticklabel style={/pgf/number format/fixed,/pgf/number format/precision=2},
    nodes near coords={\pgfmathprintnumber[fixed,fixed zerofill,precision=3]{\pgfplotspointmeta}},
    nodes near coords style={font=\scriptsize, color=axisO!80!black},
    point meta=y]
  \addplot[ybar, bar width=20pt, fill=axisOf, draw=axisO]
    coordinates {(1,0.891) (2,0.907) (3,0.913) (4,0.918) (5,0.931)};
\end{axis}
\begin{axis}[progaxis,
    ymin=0, ymax=10, ylabel={Non-referral rate (\%) (lines)}, axis y line*=right, axis x line=none,
    xticklabels={}, ytick={0,2,4,6,8,10},
    legend style={font=\scriptsize, at={(0.5,1.03)}, anchor=south, draw=black!40, legend columns=2}]
  \addplot[very thick, failred, mark=*, mark options={fill=white, draw=failred},
           nodes near coords={\pgfmathprintnumber[fixed,precision=1]{\pgfplotspointmeta}\%},
           nodes near coords style={font=\scriptsize, color=failred, anchor=north west, xshift=2pt, yshift=-1pt},
           point meta=y]
    coordinates {(1,5.8) (2,4.1) (3,3.7) (4,2.9) (5,1.4)};
  \addlegendentry{Sev-NR}
  \addplot[thick, dashed, axisR, mark=square*, mark options={fill=white, draw=axisR}]
    coordinates {(1,8.6) (2,6.2) (3,5.1) (4,4.0) (5,1.9)};
  \addlegendentry{PDR-NR}
  \draw[black!60, densely dotted, thick] (axis cs:0.4,2.0) -- (axis cs:5.6,2.0)
    node[pos=0.03, anchor=south west, font=\scriptsize, fill=white, inner sep=1pt, text=black!70] {$2\%$ threshold};
\end{axis}
\end{tikzpicture}
\caption[DOTS ablation variants on Messidor-2]{DOTS ablation variants evaluated on
Messidor-2 ($1{,}744$ images never used for training; values of \cref{tab:h3b}). Bars:
QWK (left axis); lines: Sev-NR and PDR-NR (right axis). Only the complete pipeline brings
both non-referral rates to $2\%$ or below.}
\label{fig:progressive}
\end{figure}


\section{Discussion}
\label{sec:ch4-discussion}

\subsection{Interpreting the Seven-Criterion Result}
\label{sec:interp7}

The ``$7/7$ pass'' result of \cref{tab:h3a} establishes that, on five
retrospective public benchmarks with stratified bootstrap intervals, DOTS's
point estimates clear every threshold drawn from published national screening
guidelines. It does not establish that DOTS will continue to clear these
thresholds on prospective data, nor that the specific threshold values are
themselves optimal. Clearing the bar on retrospective data is Level~1 evidence
in the hierarchy of the General Conclusion, not Level~2--4 evidence. The
chapter therefore says that DOTS \emph{satisfies} the seven criteria, not that
it is \emph{proven safe}.

\subsection{Claim--Evidence--Strength Mapping}
\label{sec:claims}

\Cref{tab:claims} maps each scientific claim of the thesis to its evidence and gives an
explicit assessment of the strength of that evidence, so that the reader can calibrate
confidence in each result. The classification is deliberately conservative: a claim is called strong only when it rests on a pre-specified test with a large effect and at least one supporting analysis.

\begin{table}[H]
\centering
\caption[Claim--evidence--strength mapping]{Claim--evidence--strength mapping. Strong: pre-specified test, large effect,
supporting analyses. Moderate: criterion met, but not with full statistical confidence
or only within a stated scope. Suggestive: consistent with the claim but not
confirmatory. Projected: extrapolation beyond the data.}
\label{tab:claims}
\footnotesize
\begin{tabular}{@{}P{3.6cm}P{9.3cm}l@{}}
\toprule
Claim & Evidence & Strength \\
\midrule
H1: dual-metric graph improves QWK & Pre-specified $t$-test ($t_4=3.82$, $p=0.009$, $d_z=1.71$); bootstrap CI excludes 0; exact permutation $p=0.031$ & Strong \\
H2: $H_0/H_1$ association with severity & Pre-specified ANOVA ($F=1308.9$, $\eta^2=0.51$, Bonferroni) & Strong \\
H3a: $7/7$ criteria met & Point-estimate pass (Sev-NR $=1.9\%$); 95\% CI $[1.4\%,\,2.4\%]$ exceeds the threshold & Moderate \\
H3b: no proper subset sufficient & Exhaustive $6/6$ enumeration; bootstrap CI of the closest case excludes 0 & Moderate \\
Super-additive interaction (\cref{ch:toporet}) & $+0.040 > +0.026$; consistent with an interaction model & Suggestive \\
Synergistic gains (\cref{ch:found}) & $7.0 > 6.3$~pp; consistent with complementary information & Suggestive \\
Cross-dataset generalisation & OOD $\Delta$QWK $\geq$ in-distribution $\Delta$QWK; consistent with stability & Suggestive \\
About $600$ fewer missed severe cases per year & Projection at the point estimate, assumed prevalence (\cref{tab:impact}) & Projected \\
\bottomrule
\end{tabular}
\end{table}

\subsection{What the Results Do Not Establish}
\label{sec:notestablished}

Scientific integrity requires explicit statements of what a thesis does not prove. Six
boundaries apply.
\begin{enumerate}
\item \textbf{No prospective clinical safety claim.} Sev-NR $=1.9\%$ was measured
retrospectively on curated, label-verified datasets, and its 95\% CI upper bound
($2.4\%$) exceeds the threshold. A prospective reader study ($N\geq200$ patients, at
least three sites) is required before any claim of clinically validated safety.
\item \textbf{No demographic equity guarantee.} None of the five benchmarks carries
per-image demographic annotations; fairness across age, sex and ethnicity has not been
evaluated.
\item \textbf{No longitudinal validity.} The population graph is static (one visit per
patient). Temporal graphs using disease progression are a research perspective.
\item \textbf{No regulatory certification.} The thesis demonstrates a regulatory-oriented
design consistent with IEC~62304, ISO~14971 and the EU AI Act. It is not a regulatory
submission.
\item \textbf{The necessity claim is scoped.} H3b shows that the six proper subsets of
DOTS fail Sev-NR $<2\%$ within the evaluated family and specification, not that no
conceivable two-axis architecture could succeed.
\item \textbf{The transfer from distributed systems is analogical.} \Cref{ch:found}
provides architecturally analogous evidence for the genericity of the GraphSAGE
structural-invariant approach, not a formal transfer guarantee.
\end{enumerate}

\subsection{Outcome of the Hypotheses}

\Cref{tab:hypsummary} brings together the outcome of the three hypotheses tested in
\cref{ch:toporet} and in this chapter.

\begin{table}[H]
\centering
\caption{Outcome of the pre-specified hypotheses.}
\label{tab:hypsummary}
\small
\begin{tabular}{@{}l P{3.6cm} P{4.2cm} l l@{}}
\toprule
Hypothesis & Pre-specified test & Result & Effect & Status \\
\midrule
H1 -- relational & Paired $t$-test, five folds & $t_4=3.82$, $p=0.009$ & $d_z=1.71$ & Confirmed \\
H2 -- structural & One-way ANOVA, Bonferroni & $F(4,4995)=1308.9$, $p<0.001$ & $\eta^2=0.51$ & Confirmed \\
H3a -- sufficiency & Seven criteria, five benchmarks & $7/7$ met at the point estimate & -- & Met$^{\ast}$ \\
H3b -- necessity & Exhaustive enumeration & $6/6$ proper subsets fail & -- & Confirmed \\
\bottomrule
\end{tabular}
\tabnote{$^{\ast}$The upper bound of the 95\% confidence interval of Sev-NR ($2.4\%$) still
exceeds the $2\%$ threshold.}
\end{table}

\section*{Conclusion}
\phantomsection\addcontentsline{toc}{section}{Conclusion}

This chapter presented DOTS and evaluated it from four complementary angles. Sufficiency:
all seven predefined criteria are met at the point estimate on five public benchmarks,
with QWK $=0.951$, Sev-NR $=1.9\%$ (95\% CI $[1.4\%,\,2.4\%]$) and $0\%$ rank violations.
Necessity: none of the six systems obtained by removing one or two axes reaches
Sev-NR $<2\%$; the closest, without topology, misses by a margin whose bootstrap interval
excludes zero. Magnitude: all primary effect sizes are large, and under explicit
prevalence assumptions the Sev-NR reduction corresponds to about $600$ fewer missed severe
cases per year in a programme of $500{,}000$ patients. Limits: the evidence is
retrospective, and the upper confidence bound of Sev-NR still crosses the threshold.

A grading system that is accurate and safe on retrospective data is not yet a medical
device. The next chapter describes how DOTS was engineered for deployment: its software
architecture, its verification, and the mapping of its design to the regulatory
requirements of medical-device software.
